
% !TeX root = resumen_papers.tex
% !TeX root = main_resumen.tex
\documentclass[11pt, a4paper]{article}

% Configuración de idioma y codificación
\usepackage[T1]{fontenc}
\usepackage[utf8]{inputenc}
% Mejorar compatibilidad con captions y hooks
\usepackage{caption}
\usepackage[spanish]{babel}
\usepackage{csquotes}
% Control float placement
\usepackage{placeins}

% Configuración de bibliografía
\usepackage[backend=biber, style=apa]{biblatex}
\addbibresource{referencias.bib}

\begin{document}

\textbf{A Multimodal Vision Foundation Model for Clinical Dermatology (PanDerm)} \parencite{yan_multimodal_2025}
Este estudio introduce PanDerm, un modelo fundacional multimodal diseñado para dermatología, pre-entrenado mediante aprendizaje auto-supervisado en más de 2 millones de imágenes de diversas modalidades
(dermatoscopia, clínica, patología). El modelo fue evaluado exhaustivamente en 28 tareas distintas, logrando un rendimiento de estado del arte y demostrando una gran capacidad de generalización incluso
con escasos datos etiquetados. En estudios comparativos con lectores humanos, PanDerm superó a los clínicos en la detección temprana de melanoma en un 10.2\% y ayudó a mejorar la precisión diagnóstica 
general de los especialistas en un 11\%. Los autores concluyen que los modelos fundacionales multimodales pueden optimizar significativamente el flujo de trabajo clínico, desde el diagnóstico diferencial
hasta el pronóstico de metástasis.

\textbf{Derm1M: A Million-scale Vision-Language Dataset Aligned with Clinical Ontology} \parencite{yan_derm1m_2025}
Este trabajo presenta Derm1M, un conjunto de datos de visión-lenguaje a gran escala que contiene un millón de pares de imagen-texto, curados y alineados con una ontología clínica estandarizada de 390
condiciones de la piel. Para garantizar la calidad de los datos provenientes de diversas fuentes de internet, se implementó un proceso de refinamiento utilizando un "Curator-LLM" que filtra el ruido y
estructura las descripciones médicas. Las pruebas realizadas demuestran que los modelos entrenados con Derm1M superan significativamente a los modelos de referencia anteriores (como PubMedCLIP) en tareas
de clasificación \textit{zero-shot} y recuperación multimodal cruzada. La conclusión principal es que la alineación rigurosa de los datos con el conocimiento clínico estructurado es más crítica para el
rendimiento del modelo que simplemente aumentar el volumen de datos sin procesar.

\textbf{Resource-efficient Medical Vision Language Model for Dermatology (SCALEMED)} \parencite{yilmaz_resource-efficient_nodate}
Este paper propone SCALEMED, un marco de trabajo para desarrollar modelos de visión-lenguaje eficientes en recursos y preservadores de la privacidad, utilizando datos de entrenamiento totalmente 
sintéticos. Los autores generaron "Dermasynth", un dataset multimodal creado mediante la orquestación de LLMs para generar descripciones clínicas y modelos de difusión para sintetizar las imágenes
correspondientes. Al entrenar modelos ligeros (como Dermatollama y DermaVLM) exclusivamente con estos datos, lograron un rendimiento competitivo en tareas de VQA (Visual Question Answering), superando
incluso a modelos generales masivos como GPT-4V en métricas específicas. El estudio concluye que los datos sintéticos de alta calidad son una alternativa viable a los datos reales restringidos, permitiendo
el desarrollo de herramientas de IA médica accesibles y privadas.

\end{document}